\documentclass{article}
\usepackage[T1]{fontenc}
\usepackage{lmodern}
\usepackage{graphicx}
\usepackage{amsmath,amssymb}
\usepackage[a4paper,margin=2cm]{geometry}
\usepackage[absolute,overlay]{textpos}
\setlength{\TPHorizModule}{1mm}
\setlength{\TPVertModule}{1mm}
\usepackage[indonesian]{babel}
\usepackage{hyperref}
\setlength{\emergencystretch}{2em}
% unit: Halaman 31

\begin{document}

% === Halaman 31 (python/native) ===

31

B. Cipher Transposisi

• Cipherteks diperoleh dengan mengubah posisi huruf di dalam plainteks. 

• Dengan kata lain, algoritma ini melakukan transpose terhadap rangkaian 

huruf di dalam plainteks.

• Nama lain untuk metode ini adalah cipher permutasi, karena transpose 

setiap karakter di dalam teks sama dengan mempermutasikan karakter-

karakter tersebut.

• Contoh cipher transposisi: Columnar Transposition Cipher, Rail Fence 

Transposition Cipher

\end{document}
